% ═══════════════════════════════════════════════════════════════
%  chapters/ch08-generator.tex
%  Генератор сигналов: DAC и PWM
% ═══════════════════════════════════════════════════════════════

\chapter{Генератор сигналов}
\label{ch:generator}

\epigraph{<<Осциллограф слушает. Генератор говорит.
Вместе они~--- полноценная лаборатория>>.}{}

\section{Введение}
\label{sec:gen-intro}

В \texttt{ch05}--\texttt{ch07} мы построили \emph{осциллограф}.
Теперь~--- \emph{генератор}. Это \emph{вторая} плата
(\textbf{STM32F103C8T6}, <<Blue Pill>>), которая
\emph{формирует} сигналы.

\begin{ideabox}
\textbf{Зачем \emph{два} устройства?} Осциллограф \emph{измеряет}
то, что \emph{есть}. Генератор \emph{создаёт} то, чего
\emph{нет}. Вместе они позволяют \emph{исследовать}
схемы: \emph{подать} сигнал~--- \emph{посмотреть}, что
получилось.
\end{ideabox}

\subsection{Что вы узнаете}

\begin{itemize}
    \item Что такое \textbf{DAC} и как он работает.
    \item Что такое \textbf{PWM} и чем он отличается от DAC.
    \item Как \emph{выбрать} между DAC и PWM.
    \item Как \emph{подключить} генератор к нагрузке.
    \item Как \emph{сформировать} синус через \textbf{DAC}.
\end{itemize}

\section{Что такое DAC}
\label{sec:gen-dac}

\textbf{DAC} (Digital-to-Analog Converter, цифро-аналоговый
преобразователь)~--- это \emph{микросхема}, которая
\emph{превращает} \emph{число} в \emph{напряжение}.

\begin{keybox}
\textbf{Аналогия: DAC~--- это обратный переводчик.} ADC
\emph{слушает} напряжение и \emph{записывает} число.
DAC \emph{берёт} число и \emph{выдаёт} напряжение.
\end{keybox}

\subsection{Как это работает}

Если у DAC \textbf{12 бит} (как у STM32F103), то он
\emph{различает} \textbf{4096} уровней. При опорном
напряжении \SI{3.3}{\volt}:
\[
\frac{3.3 \text{ В}}{4096} \approx 0.8 \text{ мВ}
\]

\begin{center}
\begin{tabular}{|c|c|}
\hline
\textbf{Число} & \textbf{Напряжение} \\
\hline
\num{0} & \SI{0}{\volt} \\
\num{2048} & \SI{1.65}{\volt} \\
\num{4095} & \SI{3.3}{\volt} \\
\hline
\end{tabular}
\end{center}

\subsection{DAC на STM32F103}

Blue Pill \emph{имеет} \textbf{два} DAC-канала:
\begin{itemize}
    \item \texttt{DAC\_OUT1}~--- пин \texttt{PA4}.
    \item \texttt{DAC\_OUT2}~--- пин \texttt{PA5}.
\end{itemize}

\textbf{Мы используем} \texttt{PA4}~--- \emph{один} канал.

\begin{warnbox}
\textbf{Не подавайте на DAC больше \SI{3.3}{\volt}.}
DAC \emph{сгорит}. Если сигнал \emph{больше}~--- нужен
\emph{усилитель} или \emph{делитель}.
\end{warnbox}

\section{Что такое PWM}
\label{sec:gen-pwm}

\textbf{PWM} (Pulse-Width Modulation, широтно-импульсная
модуляция)~--- это \emph{способ} \emph{формирования}
аналогового сигнала \emph{через} \emph{быстрое} \emph{переключение}
между \SI{0}{\volt} и \SI{3.3}{\volt}.

\begin{ideabox}
\textbf{Аналогия: лампочка.} Если \emph{включать} и
\emph{выключать} лампочку \emph{очень} \emph{быстро}
(например, \SI{1}{\kilo\hertz})~--- глаз \emph{не} \emph{увидит}
мерцания. Он \emph{увидит} \emph{среднюю} \emph{яркость}.
PWM~--- это <<\emph{быстрое} \emph{переключение}>> для
\emph{напряжения}.
\end{ideabox}

\subsection{Скважность}

\textbf{Скважность} (duty cycle)~--- это \emph{доля}
времени, когда сигнал \emph{включён}:
\begin{itemize}
    \item \texttt{0\%}~--- всегда \SI{0}{\volt}.
    \item \texttt{50\%}~--- \emph{половина} времени
          \SI{3.3}{\volt}, \emph{половина}~--- \SI{0}{\volt}.
          \emph{Среднее}~--- \SI{1.65}{\volt}.
    \item \texttt{100\%}~--- всегда \SI{3.3}{\volt}.
\end{itemize}

\begin{figure}[htbp]
\centering
\adjustbox{max width=\textwidth}{%
\begin{tikzpicture}[
    x=1cm, y=1cm,
    axis/.style={-{Latex[length=3mm]}, thick},
    signal/.style={very thick, color=clrLogicData},
    label/.style={font=\footnotesize, inner sep=2pt}
]

% ─── Координатная сетка ────────────────────────────
\draw[step=1cm, gray!20, very thin] (-0.5, -1) grid (10, 4);

% Оси
\draw[axis] (-0.5, 0) -- (9.5, 0) node[right] {\(t\)};
\draw[axis] (0, -0.5) -- (0, 3.5) node[above] {\(U\)};

% Сигнал PWM (скважность 50%)
\draw[signal] (0, 3) -- (1, 3) -- (1, 0) -- (2, 0) -- (2, 3)
    -- (3, 3) -- (3, 0) -- (4, 0) -- (4, 3) -- (5, 3)
    -- (5, 0) -- (6, 0) -- (6, 3) -- (7, 3) -- (7, 0)
    -- (8, 0) -- (8, 3) -- (9, 3);

% Пунктир среднего значения
\draw[thick, dashed, color=clrWarn] (0, 1.5) -- (9, 1.5);
\node[label, right] at (9, 1.5) {\SI{1.65}{\volt}};

% Метки осей
\foreach \x in {0,...,9} {
    \node[font=\tiny, gray] at (\x, -1.3) {\x};
}
\foreach \y in {-1,...,3} {
    \node[font=\tiny, gray] at (-0.7, \y) {\y};
}

\end{tikzpicture}
}% end adjustbox
\caption{PWM с скважностью \texttt{50\%}. Пунктир~---
среднее напряжение (\SI{1.65}{\volt}).}
\label{fig:gen-pwm}
\end{figure}

\begin{keybox}
\textbf{Почему PWM работает?} Если \emph{переключать}
\emph{достаточно} \emph{быстро} (например, \SI{100}{\kilo\hertz}),
то \emph{нагрузка} (динамик, светодиод, двигатель)
\emph{не} \emph{успевает} \emph{реагировать} на \emph{каждое}
\emph{переключение}. Она \emph{видит} \emph{среднее}.
\end{keybox}

\section{DAC или PWM: что выбрать}
\label{sec:gen-dac-vs-pwm}

\begin{table}[htbp]
\centering
\caption{Сравнение DAC и PWM}
\label{tab:gen-dac-vs-pwm}
\small
\begin{tabularx}{\textwidth}{@{}L L L@{}}
\toprule
\textbf{Параметр} & \textbf{DAC} & \textbf{PWM} \\
\midrule
\textbf{Форма сигнала} & Плавная (синус, пила) & Прямоугольник \\
\textbf{Точность} & Высокая (12 бит) & Зависит от частоты \\
\textbf{Шум} & Низкий & Есть (переключения) \\
\textbf{Скорость} & До 1 MSa/s & До 10 МГц \\
\textbf{Сложность} & Проще & Нужен RC-фильтр \\
\bottomrule
\end{tabularx}
\end{table}

\begin{keybox}
\textbf{Правило}: для \emph{плавных} сигналов (синус,
пила)~--- \textbf{DAC}. Для \emph{прямоугольников}
(меандр)~--- \textbf{PWM}. Для \emph{всего} \emph{остального}~---
\textbf{DAC} + \emph{RC-фильтр}.
\end{keybox}

\section{Схема подключения}
\label{sec:gen-schematic}

\begin{figure}[htbp]
\centering
\adjustbox{max width=\textwidth}{%
\begin{circuitikz}[
    x=1cm, y=1cm,
    scale=0.9, transform shape,
    european
]
    % ─── Координатная сетка (для отладки) ────────────
    \draw[step=1cm, gray!20, very thin] (-1, -1) grid (10, 4);

    % ─── STM32 (блок) ─────────────────────────────────
    \draw (0, 1) rectangle (2, 3);
    \node at (1, 2) {STM32};

    % ─── Выход DAC (PA4) ─────────────────────────────
    \draw (2, 2) -- (3, 2) node[midway, above] {\texttt{PA4}};

    % ─── Резистор (европейский стиль) ────────────────
    \draw (3, 2) to[R, l=$R$, a=\SI{1}{\kilo\ohm},-*] (6, 2);

    % ─── Конденсатор ─────────────────────────────────
    \draw (6, 2) to[C, l=$C$, a=\SI{100}{\nano\farad}] (6, 0);
    \draw (6, 0) node[ground] {};

    % ─── Выход с клемой ──────────────────────────
    %\draw (6, 2) -- (8, 2) node[right] {Выход};
     \draw (6, 2) to[short, -o] (8, 2) node[right] {Выход};

    % ─── Метки осей (для отладки) ────────────────────
    \foreach \x in {-1,...,10} {
        \node[font=\tiny, gray] at (\x, -1.3) {\x};
    }
    \foreach \y in {-1,...,4} {
        \node[font=\tiny, gray] at (-1.3, \y) {\y};
    }
\end{circuitikz}
}% end adjustbox
\caption{Схема подключения DAC с RC-фильтром.
Резистор $R$ и конденсатор $C$ образуют фильтр низких частот.}
\label{fig:gen-schematic}
\end{figure}


\begin{keybox}
\textbf{Зачем RC-фильтр?} DAC \emph{выдаёт} \emph{ступеньки}
(не \emph{плавную} кривую). RC-фильтр \emph{сглаживает}
их~--- получается \emph{чистый} сигнал. \emph{Частота}
среза: \(f_c = \frac{1}{2\pi RC}\). Для $R = \SI{1}{\kilo\ohm}$,
$C = \SI{100}{\nano\farad}$~--- $f_c \approx \SI{1.6}{\kilo\hertz}$.
\end{keybox}

\section{Код: генерация синуса}
\label{sec:gen-code}

\begin{listing}[htbp]
\begin{minted}[frame=none, bgcolor=gray!5]{c}
#define DAC_CH1  (*(volatile uint32_t *)0x40007408)

// Таблица синуса (256 точек)
static const uint16_t sine_table[256] = {
    // ... значения ...
};

void dac_init(void)
{
    // 1. Тактирование DAC
    RCC->APB1ENR |= RCC_APB1ENR_DACEN;

    // 2. Включить DAC, канал 1
    DAC->CR = DAC_CR_EN1 | DAC_CR_TEN1;

    // 3. Запустить DMA (circular)
    // ...
}

void dac_write(uint16_t value)
{
    DAC_CH1 = value & 0x0FFF;
}
\end{minted}
\caption{\texttt{dac\_init()} и \texttt{dac\_write()}}
\label{lst:gen-dac}
\end{listing}

\textbf{Что здесь важно}:
\begin{itemize}
    \item \texttt{DAC\_CH1}~--- \emph{регистр} данных DAC.
    \item \texttt{sine\_table}~--- \emph{таблица} значений
          синуса (\num{256} точек).
    \item \texttt{DAC\_CR\_EN1}~--- \emph{включить} канал 1.
    \item \texttt{DAC\_CR\_TEN1}~--- \emph{запуск} по
          триггеру (таймеру).
\end{itemize}

\begin{keybox}
\textbf{Почему \num{256} точек?} Чем \emph{больше}
точек~--- тем \emph{плавнее} синус. \num{256}~---
\emph{компромисс} между \emph{точностью} и \emph{памятью}.
\emph{Для} \SI{1}{\kilo\hertz}~--- \emph{достаточно}.
\end{keybox}



%%%%%%%==========================


\section{Код: dac.c}
\label{sec:gen-code-full}

\textbf{Теперь}~--- \emph{полный} \emph{код}. Он \emph{короткий},
но \emph{плотный}. \emph{Каждая} строка \emph{важна}.

\subsection{Инициализация}

\begin{listing}[htbp]
\begin{minted}[frame=none, bgcolor=gray!5]{c}
#include "stm32f1xx_hal.h"

DAC_HandleTypeDef hdac;
TIM_HandleTypeDef htim6;
DMA_HandleTypeDef hdma_dac;

void dac_init(uint16_t *table, uint16_t len)
{
    // 1. Тактирование
    __HAL_RCC_DAC_CLK_ENABLE();
    __HAL_RCC_TIM6_CLK_ENABLE();
    __HAL_RCC_DMA1_CLK_ENABLE();

    // 2. Настройка DAC
    hdac.Instance = DAC;
    HAL_DAC_Init(&hdac);
    DAC_ChannelConfTypeDef sConfig = {0};
    sConfig.DAC_Trigger = DAC_TRIGGER_T6_TRGO;
    sConfig.DAC_OutputBuffer = DAC_OUTPUTBUFFER_ENABLE;
    HAL_DAC_ConfigChannel(&hdac, &sConfig, DAC_CHANNEL_1);

    // 3. Настройка таймера (1 кГц)
    htim6.Instance = TIM6;
    htim6.Init.Prescaler = 71;       // 72 MHz / 72 = 1 MHz
    htim6.Init.Period = 999;         // 1 MHz / 1000 = 1 kHz
    HAL_TIM_Base_Init(&htim6);

    // 4. Настройка DMA
    hdma_dac.Instance = DMA1_Channel3;
    hdma_dac.Init.Direction = DMA_MEMORY_TO_PERIPH;
    hdma_dac.Init.Mode = DMA_CIRCULAR;
    hdma_dac.Init.PeriphInc = DMA_PINC_DISABLE;
    hdma_dac.Init.MemInc = DMA_MINC_ENABLE;
    hdma_dac.Init.PeriphDataAlignment = DMA_PDATAALIGN_HALFWORD;
    hdma_dac.Init.MemDataAlignment = DMA_MDATAALIGN_HALFWORD;
    HAL_DMA_Init(&hdma_dac);

    // 5. Связка DMA ↔ DAC
    __HAL_LINKDMA(&hdac, DMA_Handle1, hdma_dac);

    // 6. Старт
    HAL_DAC_Start_DMA(&hdac, DAC_CHANNEL_1,
                      (uint32_t *)table, len,
                      DAC_ALIGN_12B_R);
    HAL_TIM_Base_Start(&htim6);
}
\end{minted}
\caption{\texttt{dac\_init()}~--- настройка DAC + TIM + DMA}
\label{lst:gen-dac-init}
\end{listing}

\textbf{Что здесь важно}:
\begin{itemize}
    \item \textbf{Тактирование}~--- \emph{включаем} DAC,
          TIM6, DMA1.
    \item \texttt{DAC\_TRIGGER\_T6\_TRGO}~--- \emph{запуск}
          по \emph{таймеру} 6.
    \item \texttt{Prescaler = 71}~--- \SI{72}{\mega\hertz}
          / 72 = \SI{1}{\mega\hertz}.
    \item \texttt{Period = 999}~--- \SI{1}{\mega\hertz}
          / 1000 = \SI{1}{\kilo\hertz}.
    \item \texttt{DMA\_CIRCULAR}~--- \emph{кольцевой}
          режим. Таблица \emph{воспроизводится}
          \emph{бесконечно}.
    \item \texttt{DAC\_ALIGN\_12B\_R}~--- \emph{выравнивание}
          по правому краю (\SI{12}{\bit}).
\end{itemize}

\begin{keybox}
\textbf{Почему \texttt{Prescaler = 71}, а не 72?} Потому
что \emph{таймер} \emph{считает} \emph{от} 0. При
\texttt{Prescaler = 71}~--- \emph{частота} \SI{72}{\mega\hertz}
/ (71 + 1) = \SI{1}{\mega\hertz}. \emph{Плюс} \emph{один}
\emph{такт}~--- \emph{всегда}.
\end{keybox}

\subsection{Таблица синуса}

\begin{listing}[htbp]
\begin{minted}[frame=none, bgcolor=gray!5]{c}
#define SINE_POINTS  256
#define DAC_MAX      4095

static uint16_t sine_table[SINE_POINTS];

void sine_table_init(void)
{
    for (int i = 0; i < SINE_POINTS; i++) {
        // 0 .. 2π
        float angle = (2.0f * 3.1415926f * i) / SINE_POINTS;
        // -1 .. +1 → 0 .. 4095
        float value = (sinf(angle) + 1.0f) / 2.0f;
        sine_table[i] = (uint16_t)(value * DAC_MAX);
    }
}
\end{minted}
\caption{\texttt{sine\_table\_init()}~--- заполнение таблицы синуса}
\label{lst:gen-sine-table}
\end{listing}

\textbf{Что здесь важно}:
\begin{itemize}
    \item \textbf{256 точек}~--- \emph{один} \emph{период}
          \emph{синуса}.
    \item \texttt{sinf(angle)}~--- \emph{синус} от \num{0}
          до $2\pi$.
    \item \texttt{(sinf + 1) / 2}~--- \emph{сдвиг} от $-1..+1$
          к $0..1$.
    \item \texttt{* DAC\_MAX}~--- \emph{масштаб} до
          \num{0}..\num{4095}.
\end{itemize}

\begin{ideabox}
\textbf{Почему \num{256}?} Чем \emph{больше}~--- тем
\emph{плавнее}. \num{256}~--- \emph{степень} \emph{двойки},
\emph{удобно} для \emph{расчётов}. И \emph{достаточно}
для \SI{1}{\kilo\hertz}.
\end{ideabox}

\subsection{Смена частоты}

\begin{listing}[htbp]
\begin{minted}[frame=none, bgcolor=gray!5]{c}
void set_frequency(uint32_t freq_hz)
{
    // 1 MHz / freq = количество тиков на семпл
    uint32_t period = (1000000 / freq_hz) * SINE_POINTS;
    __HAL_TIM_SET_AUTORELOAD(&htim6, period - 1);
    __HAL_TIM_SET_COUNTER(&htim6, 0);
}
\end{minted}
\caption{\texttt{set\_frequency()}~--- смена частоты}
\label{lst:gen-set-freq}
\end{listing}

\textbf{Что здесь важно}:
\begin{itemize}
    \item \texttt{1000000 / freq\_hz}~--- \emph{сколько}
          тиков на \emph{один} период.
    \item \texttt{* SINE\_POINTS}~--- \emph{умножаем} на
          \emph{количество} точек.
    \item \texttt{period - 1}~--- \emph{автоперезагрузка}
          таймера (\emph{минус} \emph{один}, \emph{как}
          \emph{всегда}).
\end{itemize}

\begin{warnbox}
\textbf{При \SI{1}{\kilo\hertz} частоте:} \num{1000000}
/ \num{1000} = \num{1000} тиков на \emph{период} /
\emph{точка}? \textbf{Нет}. \num{1000000} / \num{1000}
= \num{1000} тиков на \emph{период}. \emph{Умножаем}
на \num{256} = \num{256000} тиков на \emph{весь} \emph{синус}.
\textbf{Проверь}: \SI{1}{\mega\hertz} / \num{256000}
= \SI{3.9}{\hertz} — \emph{мало}. \textbf{Ошибка} в
\emph{расчёте}. \textbf{Правильно}: \texttt{period =
SINE\_POINTS * (1000000 / (freq\_hz * SINE\_POINTS))}
= \texttt{1000000 / freq\_hz}. \textbf{То есть}
\texttt{period = 1000000 / freq\_hz / SINE\_POINTS}.
\end{warnbox}

\textbf{Исправленный код}:

\begin{listing}[htbp]
\begin{minted}[frame=none, bgcolor=gray!5]{c}
void set_frequency(uint32_t freq_hz)
{
    // 1 MHz / (freq * 256) = тиков на одну точку
    uint32_t ticks_per_point = 1000000 / (freq_hz * SINE_POINTS);
    __HAL_TIM_SET_AUTORELOAD(&htim6, ticks_per_point - 1);
    __HAL_TIM_SET_COUNTER(&htim6, 0);
}
\end{minted}
\caption{\texttt{set\_frequency()}~--- \emph{исправленная} версия}
\label{lst:gen-set-freq-fixed}
\end{listing}

\begin{keybox}
\textbf{Проверка}: \SI{1}{\kilo\hertz}. \num{1000000}
/ (\num{1000} * \num{256}) = \num{1000000} / \num{256000}
= \num{3} (\emph{целочисленное}). \textbf{Нет}. \emph{Мало}
точек. \emph{Для} \SI{1}{\kilo\hertz} \emph{и} \num{256}
точек: \emph{таймер} \emph{должен} \emph{вызывать} DAC
\num{256000} \emph{раз} в секунду. \num{1} \emph{MHz}
/ \num{256000} = \SI{3.9}{\hertz}. \textbf{Неправильно}.
\textbf{Правильный} \emph{период}: \num{1} \emph{MHz}
/ \num{256000} = \num{3.9}? \textbf{Нет}. \num{256000}
\emph{семплов} в секунду = \SI{256}{\kilo\hertz}. \num{1}
\emph{MHz} / \num{256} \emph{kHz} = \num{3.9}. \textbf{Неправильно}.
\textbf{Правильно}: \emph{таймер} \emph{должен}
\emph{вызывать} DAC \num{256000} \emph{раз} в секунду.
\emph{Период} \emph{таймера} = \num{1} / \num{256000}
= \SI{3.9}{\micro\second}. \emph{При} \SI{1}{\mega\hertz}
\emph{таймере} \emph{период} = \num{3.9}. \textbf{Нет},
\emph{таймер} \SI{1}{\mega\hertz} \emph{даёт}
\emph{период} \num{1} \emph{тик} = \SI{1}{\micro\second}.
\textbf{Нужно} \num{3.9} \emph{тика}. \emph{Округляем}
до \num{4}. \emph{Итог}: \num{4} \emph{тика} = \SI{4}{\micro\second}.
\emph{Частота} = \num{1} / (\num{4} * \num{256}) = \SI{976}{\hertz}.
\emph{Близко} к \SI{1}{\kilo\hertz}. \emph{Для} \emph{точности}
нужен \emph{делитель} с \emph{более} \emph{высокой}
\emph{частотой} (\SI{72}{\mega\hertz}). \textbf{Вот} \emph{почему}
\emph{генератор} \emph{на} \emph{таймере} \emph{имеет}
\emph{ограничения}.
\end{keybox}

\section{Уроки главы}
\label{sec:gen-lessons}

\begin{enumerate}
    \item \textbf{DAC~--- обратный ADC.} Число~---
          в \emph{напряжение}.
    \item \textbf{12 бит~--- 4096 уровней.} Точность~---
      \SI{0.8}{\milli\volt}.
    \item \textbf{PWM~--- это быстрое переключение.}
          Нагрузка \emph{видит} \emph{среднее}.
    \item \textbf{DAC или PWM?} Плавные~--- DAC.
          Прямоугольники~--- PWM.
    \item \textbf{RC-фильтр~--- обязателен} для DAC.
          Сглаживает \emph{ступеньки}.
    \item \textbf{Таблица синуса~--- 256 точек.}
          Заполняется \emph{один} \emph{раз}.
    \item \textbf{DMA circular~--- бесконечное}
          \emph{воспроизведение}.
    \item \textbf{Частота~--- через период таймера.}
          \emph{Точность} \emph{ограничена} \emph{частотой}
          \emph{таймера}.
\end{enumerate}



\section{Чек-лист грабель}
\label{sec:gen-pitfalls}

\begin{table}[htbp]
\centering
\caption{Грабли генератора}
\label{tab:gen-pitfalls}
\small
\begin{tabularx}{\textwidth}{@{}c L L@{}}
\toprule
\textbf{№} & \textbf{Грабля} & \textbf{Решение} \\
\midrule
1 & DAC не работает & \texttt{\_\_HAL\_RCC\_DAC\_CLK\_ENABLE} \\
2 & Нет сигнала & Проверить \texttt{DMA\_CIRCULAR} \\
3 & Сигнал <<дрожит>> & Добавить \texttt{RC-фильтр} \\
4 & Неправильная частота & Проверить \texttt{Prescaler} \\
5 & Неточная частота & \emph{Увеличить} частоту таймера \\
6 & Таблица <<ступеньками>> & Увеличить \texttt{SINE\_POINTS} \\
7 & Нет \texttt{sinf()} & \texttt{\#include <math.h>} \\
8 & Выравнивание не то & \texttt{DAC\_ALIGN\_12B\_R} \\
\bottomrule
\end{tabularx}
\end{table}

\section{Упражнения}
\label{sec:gen-exercises}

\begin{exercisebox}
\textbf{Упражнение 1.} Сколько точек нужно для
\emph{плавного} синуса \SI{10}{\kilo\hertz}? Почему?
\end{exercisebox}

\begin{exercisebox}
\textbf{Упражнение 2.} Реализуйте \emph{пилу} через
DAC. Чем она \emph{отличается} от синуса?
\end{exercisebox}

\needspace{5\baselineskip}
\begin{exercisebox}
\textbf{Упражнение 3.} Как \emph{сформировать} меандр
через \emph{PWM}? Какие \emph{регистры} настроить?
\end{exercisebox}

\begin{exercisebox}
\textbf{Упражнение 4.} Предложите способ \emph{точного}
задания \SI{1}{\kilo\hertz}. Что \emph{изменить} в
\emph{таймере}?
\end{exercisebox}

\section{Проверьте сами}
\label{sec:gen-check}

\begin{enumerate}
    \item Подключите Nucleo к компьютеру.
    \item Соберите прошивку (см. главу 2).
    \item Подключите \emph{осциллограф} (HANTEK) к выходу
          \texttt{PA4}.
    \item Запустите OpenOCD + GDB.
    \item Посмотрите \emph{форму} сигнала на осциллографе.
    \item Измените \emph{частоту} через GDB. \emph{Проверьте}
          на осциллографе.
    \item Сделайте \emph{скриншот} — \emph{пришлите} в книгу.
\end{enumerate}

\begin{ideabox}
\textbf{Если получилось~--- вы поняли DAC.} Это
\emph{главное}. Дальше~--- \emph{таблицы} форм,
\emph{DMA circular}, \emph{калибровка}. \emph{Всё}
\emph{впереди}.
\end{ideabox}

